\section{Алгоритм с изменяемой длительностью измерительного окна, равной периоду инжектируемого сигнала, и децимацией}\label{sec:N_const}

\needspace{1\baselineskip}
Рассматриваемый алгоритм представляет собой модификацию алгоритма \ref{sec:N_var} в части периода дискретизации ${T_s}$ -- период дискретизации пропорционален периоду инжектируемого сигнала ${T_z}$. Ввиду плавающего ${T_s}$ длина $N$ последовательностей ${x\left[n\right]}$, ${y\left[n\right]}$, по которым вычисляются коэффициенты разложения в выражении \eqref{eqn:DFT}, неизменна для всех значений периода инжектируемого сигнала ${T_z}$. Условие кратности периоду выполнения операций в системе ${T_d}$ остается.

Постоянство $N$ при изменении периода инжектируемого сигнала ${T_z}$ обеспечивается децимацией (прореживанием) последовательностей ${x'[n]}$, ${y'[n]}$ с коэффициентом децимации ${M^{(k)} \in \mathbb{N}}$:
\begin{equation*}
    \begin{split}
        &x^{(k)}[n] = x'^{(k)}\left[n \mathord{\cdot} M^{(k)}\right] = x^{(k)}_c \left(n \mathord{\cdot} M^{(k)} \mathord{\cdot} T_d \right), \\
        &y^{(k)}[n] = y'^{(k)}\left[n \mathord{\cdot} M^{(k)}\right] = y^{(k)}_c \left(n \mathord{\cdot} M^{(k)} \mathord{\cdot} T_d \right),
    \end{split}
\end{equation*}
Переменными ${x'[n]}$, ${y'[n]}$ обозначены последовательности, формируемые в результате дискретизации откликов ${x_c(t)}$, ${y_c(t)}$ непосредственно с периодом выполнения операций в системе ${T_d}$.

Таким образом, в серии из $K$ измерений для каждого значения периода инжектируемого сигнала ${T^{(k)}_z}$, ${k \in \{1, \dots, K\}}$ (рисунок~\ref{fig:N_const}), выполняется:
\begin{equation*}
    T^{(k)}_h = T^{(k)}_z = N^{(k)} \mathord{\cdot} \underbrace{M^{(k)} \mathord{\cdot} T_d}_{T^{(k)}_s}.
\end{equation*}

\begin{figure}[H]
    \centering
    \includegraphics [width=1.00\linewidth] {frequencyResponse/N_const.pdf}

    \caption{Параметры в алгоритме с изменяемой длительностью измерительного окна, равной периоду инжектируемого сигнала, и децимацией}
    \label{fig:N_const}
\end{figure}

В связи с тем, что \textbf{так же, как и в алгоритме \ref{sec:N_var}}, длительность измерительного окна ${T^{(k)}_h}$ равна периоду инжектируемого сигнала ${T^{(k)}_z}$, составляющей с частотой ${f^{(k)}_z}$ соответствуют коэффициенты ${X^{(k)}\left[l\right]}$, ${Y^{(k)}\left[l\right]}$ всегда при ${l = 1}$, т.е. коэффициент передачи динамического звена на частоте ${f^{(k)}_z}$ определяется выражением \eqref{eqn:complexFreqResp} при ${L^{(k)} = 1}$. \textbf{Принципиальное отличие от алгоритма \ref{sec:N_var}} связано с потенциально меньшей длиной $N$ массивов-откликов вследствие децимации.

\needspace{1\baselineskip}
Для изложенного алгоритма измерения характерно следующее:
\begin{enumerate}
    \item Максимально возможный период инжектируемого сигнала ${T_z}$ соответствует максимальному значению измерительного окна ${T_h}$, но уже не ограничен параметром ${N_\text{max}}$, т.к. длина $N$ каждого массива-отклика ${x\left[n\right]}$, ${y\left[n\right]}$ неизменна (не зависит от периода инжектируемого сигнала ${T_z}$).
    \item Шаг приращения периода инжектируемого сигнала ${T_z}$ (разрешающая способность) в $N$ раз больше по сравнению с алгоритмом \ref{sec:N_var} (${\Delta T_z = N \mathord{\cdot} T_d}$).
    \item Грубый шаг приращения периода обусловливает еще более неравномерное распределение частот. Например, в диапазоне частот ${f \in \left[{f_d / 100}; {f_d / 10} \right]}$ при ${N = 10}$ частота инжектируемого сигнала ${f_z}$ может принимать значения всего из 10 возможных: \\
    ${f_z = \left({f_d / 100}, {f_d / 90}, ..., {f_d / 10} \right)}$.
\end{enumerate}
